% Seis signos funcionales nuevos. Propuestas gráficas pendientes de prueba manual.
\KalaRegistrar{K17}{Alternativa}{Operador}
  {Disyunción inclusiva de dos proposiciones: al menos una es verdadera; pueden serlo ambas.}{2}{%
  \draw (14,18) -- (30,32) -- (14,46);
  \draw (50,18) -- (34,32) -- (50,46);
}
\KalaRegistrar{K18}{Posibilidad}{Operador}
  {Con un índice de modelo, afirma que el cuerpo se cumple en algún estado accesible. Fragmento modal proposicional.}{2}{%
  \draw (44,14) -- (20,14) .. controls (12,14) and (12,18) .. (12,26) -- (12,42) .. controls (12,50) and (16,50) .. (24,50) -- (44,50);
  \draw (32,23) -- (43,32) -- (32,41);
}
\KalaRegistrar{K19}{Pregunta}{Marcador}
  {Pregunta si la proposición agrupada es verdadera. Solo aparece en el nivel superior; no afirma el contenido.}{1}{%
  \draw (14,43) .. controls (14,16) and (45,14) .. (49,30) -- (33,40) -- (49,50);
}
\KalaRegistrar{K20}{Selección}{Constructor}
  {Construye el conjunto de elementos de un dominio que cumplen una condición. Vincula una variable en el cuerpo.}{2}{%
  \draw (14,16) -- (50,16) -- (35,34) -- (35,49);
  \draw (22,25) -- (42,25);
}
\KalaRegistrar{K21}{Secuencia}{Programa}
  {Ejecuta el primer programa y después el segundo. El segundo comienza solo si el primero termina.}{2}{%
  \draw (14,45) -- (14,32) -- (31,32) -- (31,19) -- (49,19);
  \draw (41,11) -- (49,19) -- (41,27);
}
\KalaRegistrar{K22}{Mientras}{Programa}
  {Comprueba una condición antes de cada iteración y ejecuta el cuerpo mientras sea verdadera. Puede no terminar.}{1}{%
  % Arco y barra en un solo trayecto; se elimina el segmento superpuesto.
  \draw (46,20) .. controls (20,5) and (9,33) .. (22,46) .. controls (34,58) and (54,41) .. (47,30) -- (37,30);
}

% Alias conservados para los comandos del cuaderno de comparación.
\KalaRegistrar{V06}{Colección B}{Aceptada}
  {Forma simplificada aceptada de K06. Alias gráfico de la forma habitual.}{2}{%
  \csname KalaDibujoK06\endcsname
}
\KalaRegistrar{V16}{Fuente B}{Aceptada}
  {Forma simplificada aceptada de K16. Alias gráfico de la forma habitual.}{2}{%
  \csname KalaDibujoK16\endcsname
}
